\documentclass[10pt,a4paper]{amsart}
\usepackage[margin=19mm]{geometry}
\usepackage{amsmath,amssymb,mathtools}
\usepackage[scr=rsfs,cal=euler]{mathalfa}
\usepackage{xfrac}
\usepackage[unicode,hidelinks,bookmarks=false]{hyperref}
\newcommand{\ie}{i.e.\ }
\newcommand{\cf}{cf.\ }
\newcommand{\resp}{respectively\ }
\newcommand{\Subsection}{\subsection}
\providecommand{\nobreakdash}{\nobreak\hbox{-}}
\setlength{\emergencystretch}{2em}
\multlinegap0pt
\usepackage{cleveref}
\newtheorem{prop}{Proposition}
\newcommand*\diff{\mathop{}\!\mathrm{d}}
\newcommand{\N}{\mathbb{N}}
\newcommand{\Z}{\mathbb{Z}}
\newcommand{\R}{\mathbb{R}}
\title[Dense sets excluding nested finite-sum windows]{Dense sets excluding nested finite-sum windows}
\author[Hern\'andez, Kousek and Radi\'c]{Felipe Hern\'andez, Ioannis Kousek, Trist\'an Radi\'c}
\date{Source: arXiv:2510.18788v1 (CC BY 4.0)}
\begin{document}
\maketitle
\noindent\emph{Source.} On density analogs of Hindman's finite sums theorem,
arXiv:2510.18788v1. The fixed-version article licence is CC BY 4.0;
the source and licence evidence accompany this editable extraction.
\par\smallskip
\noindent\emph{Scope.} One counted proposition, source lines 438--441,
and its complete author proof, lines 460--484, are reproduced unchanged.
The following source conventions are uncounted context.
\par\smallskip
By $\Z$, $\N_0$, $\N$ and $\R$ we denote, respectively, the integers,
nonnegative integers, positive integers (natural numbers), and the reals. 

\par\smallskip
\noindent The source defines the upper Banach density of $A\subset\N$ by
\begin{equation*}
    \diff^{*}(A)=\limsup_{N-M\to \infty}\frac{|A\cap \{M,M+1,\ldots,N-1\}|}{N-M},
\end{equation*}

\noindent\emph{Domain note.} In the proposition, $k_i,t_i\in\N$ and
the $B_i$ are infinite subsets of $\N$, as in the article's natural-number
setting (abstract, lines 244--245; question, line 281).
Thus $p_n\N=\{p_n,2p_n,3p_n,\ldots\}$.
\par\medskip
\begin{prop}\label{1.2 is optimal (ii)}
Given any $\epsilon>0$ there exists a set $A\subset \N$ with $\diff^*(A)>1-\epsilon$ such that for any sequence $(k_i)_{i \in \N}$, for any nested sequence $B_1\supset B_2 \supset \ldots$, of infinite sets, and any sequence $(t_i)_{i\in \N}$, it does not hold that
$$\bigg\{\sum_{n\in F}n: F\subset B_i\ \text{with}\ k_i\leq |F|\leq k_i+i \bigg\}\subset A-t_i,\ \text{for every}\ i\in \N.$$
\end{prop}

\begin{proof}[Proof of \cref{1.2 is optimal (ii)}]
Let $\epsilon>0$ be fixed and consider an increasing sequence of prime numbers $(p_n)_{n\in \N}$ such that $\sum_{n=1}^{\infty} n/p_n <\epsilon.$ Then, we define the set 
\begin{equation} \label{eq set counterex}
    A=\N \setminus \bigcup_{n=1}^{\infty} \left( \bigcup_{j=0}^{n-1}p_n\N + j \right).
\end{equation}
By the choice of $(p_n)$ it readily follows that $\diff^*(A)>1-\epsilon$. Now, assume there exists an infinite sequence of sets $B_1 \supset B_2 \supset B_3 \supset \dots$ and a sequence of shifts $(t_i)_{i\in \N}$ such that for each $i\in \N$ it holds that
$$ \bigg\{\sum_{n\in F}n: F\subset B_i\ \text{with}\ k_i\leq |F|\leq k_i+i \bigg\} \subset A-t_i.$$
Notice that, if $r>i$, the fact that $B_{r} \subset B_i$ implies
$$\bigg\{\sum_{n\in F}n: F\subset B_r\ \text{with}\ k_i\leq |F|\leq k_i+i \bigg\} \subset A-t_i.$$
Fix $i,m,n,r\in \N$ with $t_{i}<n<p_n<r<m$ and $i<r<m$ so that $B_m \subset B_{r} \subset B_{i}$.

Assume first that there is an infinite set $\tilde{B}_m\subset B_m \cap  p_n\N$. In particular, $$\bigg\{\sum_{n\in F}n: F\subset \tilde{B}_m\ \text{with}\  |F|= k_i \bigg\} \subset p_n \N.$$  Then, by assumption and the above 
discussion, it must be the case that 
$$\bigg\{\sum_{n\in F}n: F\subset \tilde{B}_m\ \text{with}\  |F|= k_i \bigg\}+t_{i} \subset A.$$
However, $t_{i}<n$ and thus 
$$\bigg\{\sum_{n\in F}n: F\subset \tilde{B}_m\ \text{with}\  |F|= k_i \bigg\}+t_{i}  \subset \bigcup_{j=0}^{n-1} p_n\N + j \subset \N\setminus A,$$
which is a contradiction. 

Therefore, by the pigeonhole principle, we may assume there exists $c_n\in \{1,\ldots,p_n-1\}$ and an infinite set $\tilde{B}_{m} \subset B_m \cap (p_n\N+c_n)$. As before, the preceding discussion implies that
$$\bigg\{\sum_{n\in F}n: F\subset \tilde{B}_m\ \text{with}\ k_r\leq |F|\leq k_r+r \bigg\} \subset A-t_{r}.$$
Now, $p_n$ is a prime and $c_n \not\equiv 0 \pmod {p_n}$, hence $c_n$ is a generator in $\Z_{p_n}$. As $r>p_n$ this means that there exists some $d\in \{k_{r},\ldots,k_{r}+r\}$ such that $dc_n+t_{r} \equiv 0 \pmod{p_n}.$ 
We have thus established that 
$$\bigg\{\sum_{n\in F}n: F\subset \tilde{B}_m\ \text{with}\ |F|=d \bigg\}+t_r \subset A \cap \left(p_n\N + d c_n+t_{r}\right) \subset A \cap p_n\N = \emptyset,$$
which is again a contradiction. 
\end{proof}

\end{document}
